\section{OBJECTIVE AND SOURCE REVIEW}
\textbf{Erd\H{o}s \#295; independent attempt, 2026-09-23.}
Let $k(N)$ be the least number of distinct reciprocals with denominators at least $N$ summing to one. The target is $k(N)-(e-1)N\to\infty$. I read all of \texttt{295.tex}; its harmonic threshold misses arithmetic cancellation. I test the simplest additional cancellation obstruction and identify why it does not by itself sum over primes.
\section{WORK: MAXIMAL PRIME-POWER MULTIPLICITY}
In any exact identity $\sum_{i=1}^k1/n_i=1$, for each prime $p$ dividing some denominator, the largest value of $v_p(n_i)$ must occur at least twice. To prove this, multiply the identity by $L=\operatorname{lcm}(n_1,\ldots,n_k)$ and reduce modulo $p$. If just one denominator attains the largest valuation, its contribution $L/n_i$ is nonzero modulo $p$ and every other contribution is zero, whereas $L\equiv0\pmod p$. This is impossible.
Consequently, if all denominators lie in $[N,B]$ and $p$ is a prime in $(B/2,B]$, the denominator $p$ must be omitted: it is the only available multiple of $p$. More generally, any prime power which appears at its highest valuation in just one available denominator excludes that denominator from an exact reciprocal identity unless another matching valuation is added outside the range.
\section{ADVERSARIAL VERIFICATION AND REMAINING GAP}
The condition concerns the maximum valuation, not merely the number of multiples of $p$; adding a multiple of $p^2$ can introduce a new unique maximum and fail to repair the obstruction. The original problem permits arbitrarily large denominators. A single large denominator can contain many primes simultaneously, so one cannot charge a separate added denominator for each forbidden large prime. Such double counting would be an invalid proof of the desired unbounded surplus. The argument therefore improves a finite search filter and identifies a real arithmetic constraint, but does not bridge the harmonic-threshold barrier. The source's exact small-$N$ optimality computations are not independently asserted here.
\section{SOURCES CHECKED}
Complete local source: \texttt{gpt\_pro\_5.4/295.tex}; indexed official entry \url{https://www.erdosproblems.com/295}, checked on 2026-09-23. The valuation proof is elementary and no literature novelty is claimed.
\section{FINAL}
\textbf{UNRESOLVED.} A prime-power multiplicity obstruction is proved; controlling how very large denominators satisfy many constraints remains the gap.
\section{COMPLETION ESTIMATE}
Conservative subjective progress toward divergence of the surplus.
\textbf{COMPLETION: 8\%}
